\chapter{核心方程、模型变种、实现路线与实验适用性}
\label{ch:core-guide}
本章可以先于推导章节阅读。它回答两组问题：第一，双流体与Hall MHD究竟抓住哪几条核心方程，各种模型改变了什么，计算机怎样求解；第二，实验实际看见什么，Hall MHD在什么条件下适用，哪些问题已经取得有条件的解答。完整方程与推导仍在第\ref{ch:models}章，本章不重复长篇计算。

\section{最核心的方程是什么}
\subsection{双流体：分别计算离子和电子怎样运动}
对磁重联，最能体现双流体区别的，是\textbf{两套独立的动量方程及其共同电磁场耦合}：
\begin{align}
 m_in_i(\partial_t+\vel_i\cdot\grad)\vel_i
 &=+en_i(\E+\vel_i\times\B)-\grad\cdot\mathsf P_i+\vct R_i,\\
 m_en_e(\partial_t+\vel_e\cdot\grad)\vel_e
 &=-en_e(\E+\vel_e\times\B)-\grad\cdot\mathsf P_e+\vct R_e.
\end{align}
左边是各自惯性；右边依次是电磁力、压力力和碰撞交换。离子与电子质量、电荷符号不同，所以可以有不同的流速、温度、压力结构和响应尺度。二者通过
$\J=e(n_i\vel_i-n_e\vel_e)$及$\rho_q=e(n_i-n_e)$共同决定场，再由场反作用于两种流体。

这两式是\emph{识别模型的核心}，不是完整求解清单。还需各自连续性、能量或压力演化、热流闭合及场方程。完整展开见第\ref{sec:full-twofluid}节；可以直接照写的五矩闭合见第\ref{sec:full-five-moment}节，十矩见第\ref{sec:full-ten-moment}节。全Maxwell双流体通常推进$\E,\B$，电子惯性自然在电子动量中；不必额外叠加一套独立欧姆定律。\citep{Wang2015}

\subsection{Hall MHD：用总动量配合保留相对漂移的欧姆定律}
在本报告的标准无电子惯性、标量压力Hall模型中，
\begin{equation}
 \boxed{\E+\vel\times\B
   =\frac{\J\times\B}{en}-\frac{\grad p_e}{en}+\eta\J,}
 \qquad \vel\simeq\vel_i,\quad \vel_e=\vel-\frac{\J}{en}.
 \label{eq:guide-core-hall}
\end{equation}
识别Hall MHD最直接的项是$\J\times\B/(en)$。它表示离子流与电子流不同，因而用离子速度描述磁场输运时需要修正。
磁场的实际推进仍遵守Faraday定律：
\begin{equation}
 \partial_t\B=\grad\times
 \left(\vel\times\B-\frac{\J\times\B}{en}
                    +\frac{\grad p_e}{en}-\eta\J\right),
 \qquad \J=\frac{\grad\times\B}{\mu_0}.
\end{equation}
再配总连续性、总动量及热力学闭合，才有完整Hall系统，见式\eqref{eq:closed-hall-system}。
在光滑、无电阻、正压电子等条件下，磁通随电子流冻结；因此Hall项改变输运和排气结构，不单独完成电子破冻结，也不直接提供$\eta J^2$式正定加热。

\subsection{二者的关系，用四个判断记住}
\begin{enumerate}
\item 双流体先有两个物种；Hall MHD是对这套方程施加低频、准中性、电子惯性可忽略及指定闭合后的约化。
\item 改用质心速度和电流描述两物种是变量变换；真正减少物理自由度的是后续近似。
\item Hall MHD仍能重建$\vel_e$，也可以独立演化$p_i,p_e$；“两温度”并不自动等于“保留完整两物种动量”。
\item 五矩双流体已经包含Hall相对漂移和电子惯性；十矩还保留各向异性与非对角压力，但其准确性仍取决于热流闭合。
\end{enumerate}
每步代数及删项条件见第\ref{sec:twofluid-to-hall}节。论文中的名称有时较宽泛，应以实际方程为准。

\section{其他模型各改变了哪一部分}
不能把所有模型排成同一条“都是Hall MHD变种”的直线。至少要区分三个维度：场方程是否准中性、电子惯性是否保留、压力/热流用什么闭合。下面按改变的位置归类。
\begin{longtable}{P{27mm}P{61mm}P{60mm}}
\toprule 名称 & 相对共同方程所作的主要改变 & 保留或丢失的能力\\\midrule\endfirsthead
\toprule 名称 & 相对共同方程所作的主要改变 & 保留或丢失的能力\\\midrule\endhead
五矩双流体 & 每物种密度、三分量动量及标量能量；指定热流 & 保留两物种惯性；不描述完整压力各向异性。全Maxwell或准中性是另一个选择。\\
十矩双流体 & 将标量能量换成六分量压力演化；闭合三阶矩 & 可保留非对角压力；热流、相混合及不稳定模的准确性不自动保证。\\
惯性扩展MHD & 用质心/电流变量，保留电子惯性，并一致保留相应能量与应力 & 可进入电子惯性尺度；不是简单往任意Hall代码加一个$\partial_t\J$就等价。\\
电阻/超电阻Hall & 在Hall电子力平衡中加入电阻或高阶非理想项 & 给电子破冻结或小尺度正则化；系数、作用变量和热收支必须说明。\\
CGL、各向异性状态方程 & 改压力闭合，保留平行与垂直差异 & 可改变磁张力、层长与场结构；不能一般替代非旋转对称电子压力。\\
EMHD & 在目标时间和尺度上固定离子响应 & 研究电子尺度电流和磁场；不能同时预测自洽离子喷流与全局离子反馈。\\
约化双流体/陀螺流体 & 进一步施加强导场、低频、各向异性尺度或回旋平均 & 降低维数并可保留FLR；离开相应排序时不能直接推广。\\
电阻/理想MHD & 进一步忽略Hall相对漂移对宏观输运的影响；理想者再忽略非理想感应项 & 面向大尺度整体运动；不显式分辨离子、电子解耦。\\
混合动理学/PIC & 将离子或两物种改由分布函数/粒子描述 & 是不同层次的动理学描述；可作为流体对照，不应计成另一种纯Hall流体闭合。\\
\bottomrule
\end{longtable}
“闭合改变”和“惯性删项”可以组合，但组合后仍应检查守恒与适用尺度。
例如F611的动理学离子加各向异性流体电子，是混合模型；其长电子层与相近总率是闭合对照证据，不能据此声称任意十矩或Hall模型都已经复现。\citep{F611,Wang2015,Andres2016}

\section{具体怎样实现：从连续方程到一次时间步}
\subsection{两种算法共享的准备}
先写清初值、边界、归一化和模型系数，再选网格与时间推进。重联常用Harris片，但应逐物种检查力平衡；只满足总压平衡，未必满足离子与电子各自方程。磁扰动可用矢势构造，以保持初始$\grad\cdot\B=0$。
质量比、光速比、热流参数、电阻及数值滤波应独立记录。下面是求解流程，实际离散公式、参数和原文差异在第2章；这些流程不表示本报告已经运行了相应模拟。\citep{Birn2001,Wang2015}

\subsection{全Maxwell五矩/十矩双流体：推进物种状态与场}
每个网格单元存五矩变量$(n_s,m_sn_s\vel_s,\mathcal E_s)$，或十矩变量$(n_s,m_sn_s\vel_s,\mathcal P_{s,ab})$；后者若用原始二阶矩，需要减去$m_sn_su_{s,a}u_{s,b}$恢复中心压力。场可按所选离散存于单元、面或边。一次时间步通常包含：
\begin{enumerate}
\item 从守恒量恢复速度、标量压力或压力张量，检查密度、内能及压力特征值。
\item 用有限体积、DG等方法计算两物种各自的输运通量；用离散Maxwell算子推进场的旋度部分。
\item 耦合更新电磁源：物种受到$\pm en_s\E$及磁旋转，$\E$同时受到两物种电流反馈。不能分别更新后不检查相互交换的功和电荷。
\item 加入已选热流、碰撞或松弛算子；不同算子可分裂或共同隐式求解，但应检查分裂误差与刚性极限。
\item 检查电荷连续性、Gauss约束、磁散度及总能量含边界通量的残差，再接受时间步。
\end{enumerate}
显式全Maxwell传播通常有$\Delta t\lesssim C\Delta x/c$；显式电磁源还受等离子频率、回旋频率限制。局地隐式源求解可缓解部分刚性，但不自动解除剩余光波传播的限制，也不代替精度收敛。
Wang等的有限体积路线采用局地隐式源与压力保护；不同论文中的DG、全隐式方案是数值实现选择，不是新的物理闭合。\citep{Wang2015,Hakim2008}

\subsection{无电子惯性Hall MHD：由磁场求电流，再求电场}
存储$n,\rho\vel,\B$及选定的一个或两个热力学变量；没有独立的电子动量推进。一个典型步骤是：
\begin{enumerate}
\item 由$\B$的离散旋度求$\J$；将$\J,\B,n$放在相容的位置计算$\J\times\B/(en)$，并重建$\vel_e=\vel-\J/(en)$。
\item 按式\eqref{eq:guide-core-hall}组装$\E$，包括压力、指定的非理想项及其边界处理。
\item 用同一套边电场更新面磁通。若采用约束输运（CT），相邻边贡献在散度中抵消，使磁散度保持。
\item 推进总质量、总动量及热力学；电子压力按电子流输运，电阻热和人为耗散要进入相容的能量账。
\item 测量欧姆定律残差、净重联磁通、电子/离子流速和边界能流；重新检查时间步及正性。
\end{enumerate}
Hall项含一次磁场导数，外面的Faraday旋度再加一次导数，产生色散刚性。典型显式限制为
$\Delta t\lesssim C_H\Delta x^2/(\VA d_i)$，也要满足对流和电阻限制。隐式Hall推进、子循环或IMEX可以改变稳定性限制，仍需时间步收敛；Hall项不能当成普通正定扩散随意离散。\citep{Arnold2008,Toth2008}

\subsection{怎样判断实现正确，而不只看图像好看}
至少依次检查均匀平衡、磁旋转、whistler色散、线性撕裂及一个非线性重联基准。网格、时间步与闭合/耗散系数应分别扫描。
重联中最关键的额外诊断是：哪一项实际支撑X点电场？若主要来源是未记账的数值残差，即使出现了四极场和磁岛，也还没有证明模型中的破冻结机制。三维还需场线连接或适用的平行电场积分，不能只用一张二维切片。

\section{实验看到了什么：哪些支持Hall物理，哪些仍未验证}
下面汇总已核读原文。实验出现Hall特征，通常验证的是某一物理环节；不能自动认证整套流体闭合及所有尺度的预测。
\begin{longtable}{P{27mm}P{62mm}P{59mm}}
\toprule 实验/文献 & 实际现象及其支持的环节 & 尚不能据此认定什么\\\midrule\endfirsthead
\toprule 实验/文献 & 实际现象及其支持的环节 & 尚不能据此认定什么\\\midrule\endhead
MRX四极场，Ren2005 & 探针测得面外四极磁场，横向尺度与$d_i$同阶；支持解耦电流系统 & 未闭合X点全部欧姆预算，不能证明Hall项直接断线。\\
MRX双尺度层，L119 & 由磁场及密度反演的电子通道窄于离子通道；电子层约$8d_e$，与部分二维计算有差距 & 尺度趋势吻合不等于绝对层宽、排出速度与压力支撑已解释。\\
MRX导引场，GuideExp2012、L193 & 导引场增强伴随率降低、面内结构改变及下游磁压堆积 & 局地Hall效应与装置边界背压未被独立隔离。\\
MRX离子增能，Yoo2013 & 面内电势、离子流与光谱联测；先加速、下游流速趋饱和后温升 & 电功不等于不可逆热化率；碰撞、再磁化及波粒机制仍需区分。\\
MRX碰撞磁岛，E448 & 亚离子电阻层存在Hall四极与低$S$磁岛链 & 不能将无碰撞或单流体的临界阈值直接搬用；具体触发机制仍待检验。\\
MAST合并，Stanier2013 & Hall模拟结合环形压缩及真实诊断弦，可解释密度双峰形态与先后演化 & 峰值和时间尺度未逐时定量复现；等温比和增强耗散限制热分配判断。\\
\bottomrule
\end{longtable}
\citep{Ren2005,L119,GuideExp2012,L193,Yoo2013,E448,Stanier2013}

还应保留能提醒模型扩展的实验：MRX中性密度增加时Hall形态可近似保留而离子喷流变慢，需要加入中性摩擦；MAIZE的四极密度目前只支持待进一步验证的Hall解释，配套三维主计算是电阻MHD；MMS冷/热离子可在同一处分处不同磁化状态，需要更细的物种划分。这些都限制了“见到某图样就认定单离子Hall模型足够”的推断。\citep{L190,F703,L227}

\section{Hall MHD的适用条件怎样判断}
应针对目标区域及目标诊断检查，而不是只给整套装置贴上“Hall”标签。
\begin{enumerate}
\item \textbf{低频、准中性。}目标尺度通常满足$L\gg\lambda_D$，不以电荷分离和电子等离子振荡为研究对象；忽略位移电流还要检查相应低速电磁排序。鞘层、光波或高频电荷分离不属于标准Hall模型的描述范围。
\item \textbf{Hall修正对目标过程重要。}直接检查$|\J|/(enU)$及欧姆预算；在典型Alfvén排序下约为$d_i/L$。宏观盒子很大不妨碍局地薄层进入Hall区。强导场时声回旋半径与有限离子回旋尺度也可能控制解耦，不能仅用$L=d_i$作普适边界。
\item \textbf{被删的电子惯性确实较小。}检查$(m_e/e)D_e\vel_e$相对实际保留项的量级；$d_e/\delta$是常用提示，但不能替代频率、流速和几何检查。若电子扩散区由惯性主导，应采用相应惯性模型。
\item \textbf{压力和热流闭合足以回答所问问题。}标量Hall适合部分离子尺度场和平均流结构；若目标是非对角压力、强各向异性、束流或速度空间热化，需张量/动力学信息或经验证的闭合。
\item \textbf{物种、碰撞和边界正确。}Hall MHD不必严格无碰撞；电子仍可有电阻，同时离子、电子出现解耦。多离子、中性拖曳、辐射电离或装置背压重要时，须把这些过程纳入，而不是靠调大电阻统包。
\item \textbf{电子破冻结的处理可核查。}研究重联必须明确由电阻、电子惯性、压力张量或其他已说明项支撑。仅有无耗散正压Hall输运不足以提供一般的电子断线机制。
\end{enumerate}
因此“Hall MHD可用”应表述为：在指定尺度、闭合与边界下，可用来预测哪些量；它不是整场模拟从离子尺度到电子热化都有效的保证。第1章提供删项估计，第2章提供数值诊断，第5章提供实验测量细节。

\section{已解决到哪一步：逐项对应原始论文}
\label{sec:guide-solved}
以下每项给出作者引文、年份和稳定文献键。点击引文可进入参考文献，再由“全文核读档案”进入第7章；用文献键搜索HTML索引可直接打开本地原文及图号记录。
这里的“已解决”指在明确条件下建立了可检验机制或获得了可靠结果，不表示所有几何和参数下均已解决。

\begin{longtable}{P{29mm}P{63mm}P{56mm}}
\toprule 问题 & 结果与具体论文 & 证据边界\\\midrule\endfirsthead
\toprule 问题 & 结果与具体论文 & 证据边界\\\midrule\endhead
为何必须区分离子与电子流 &
\citet{Ren2005}（2005，Ren2005）图2、4测到MRX四极磁场及碰撞性关联，支持解耦电流系统；\citet{Birn2001}（2001，Birn2001）提供相关跨模型基准。 &
Ren2005的电子流及Hall电场含反演/估计，未闭合X点欧姆预算；不能写成Hall项直接断线已获验证。\\
已启动重联能否得到快演化 &
\citet{Birn2001}（2001，Birn2001）图1中Hall、混合与PIC通量曲线接近，建立了有限种子二维非线性阶段的基准。 &
强种子刻意跳过线性触发；电子层和喷流宽度仍不同。该文不证明自然启动或普适0.1。\\
同一参数下为何可有快慢两支 &
\citet{Cassak2005}（2005，Cassak2005）图3、4给出不同历史及迟滞；\citet{Cassak2010}（2010，Cassak2010）进一步在关闭电子惯性的Hall算例中保留双稳态。 &
依赖具体片层、边界和耗散；2010现象学模型未自洽选定层长，不能给出全部环境的触发阈值。\\
MAST双峰密度怎样形成 &
\citet{Stanier2013}（2013，Stanier2013）将Hall密度响应、环形内侧压缩和实验诊断弦结合，解释双峰及外峰随合并消失的顺序。 &
是形态和顺序的解释；峰值、时钟未逐时定量匹配，且采用固定温比和增强耗散。\\
X线为何单向或双向扩展 &
\citet{Shepherd2012}（2012，Shepherd2012）图3--6显示弱导场主要向电子流动方向扩展，强导场转为双向近导场Alfvén传播。 &
只检验了电子承担初始全部电流等条件；阈值和速度有诊断误差，后期斜模破坏简单图像。\\
流体能否改善电子层描述 &
\citet{Wang2015}（2015，Wang2015）图4、8--10显示十矩可改善部分局地电子流和压力结构；\citet{F611}（2016，F611）图1--7用各向异性流体电子再现长电流层。 &
Wang2015未重现完整时史及所有压力峰值；F611是混合动理学离子模型。这些不是标准标量Hall已完全解决电子层的证据。\\
\bottomrule
\end{longtable}

例如第一行的实验验证、第二行的基准一致和最后一行的闭合改善，属于不同类型的进展。不能把它们叠加成“某一套Hall代码已经同时完成全部验证”。

\section{未解决什么：原文缺口与本报告判断分别标注}
\label{sec:guide-unsolved}
以下问题均对应已核读论文。本节所说“未解决”限于这些论文提供的证据，不能据此断言所有后续文献都没有答案。“原文缺口”指文中明确保留或结果直接显示的不足；“本报告综合判断”是跨论文比较后提出的研究问题，不冒充作者原话。

\subsection{电子扩散区的绝对层宽、速度和欧姆预算}
\textbf{对应文章：}\citet{L119}（2008，L119，
\emph{Experimental study of the Hall effect and electron diffusion region during magnetic reconnection in a laboratory plasma}）；
\citet{Wang2015}（2015，Wang2015）。

\textbf{原文缺口：}L119图13的电场预算没有测全压力张量与涨落项；图16--18的电子层较所比较二维模型更宽、排出更慢，经验摩擦估计仍有差距。
Wang2015图4的重联时史不同；图8--10在等磁通、不同时间比较时，部分非对角压力峰值及下游欧姆项仍不匹配。
\textbf{本报告综合判断：}同为电子尺度的趋势一致，尚未构成对层厚、速度和全部非理想项的同时验证；下一步应在相同诊断定义和仪器平均下比较完整预算。

\subsection{同一压力/热流闭合能否跨算例迁移}
\textbf{对应文章：}\citet{Ng2020}（2020，Ng2020，
\emph{An improved ten-moment closure for reconnection and instabilities}）；
\citet{F611}（2016，F611）；
\citet{Closure2026}（2026，Closure2026）。

\textbf{原文缺口：}Ng2020图1--7及附录A显示，梯度闭合可改善岛合并和部分不稳定性，电子各向异性及跨场热流仍有误差；F611的Harris实现需要二维磁通标签及参考密度修正；Closure2026恢复了部分三维不稳定性，却仍有增长时序、各向异性位置与范围的差异。
\textbf{本报告综合判断：}尚未建立一套无需逐算例调整、又能同时预测率、层结构、热流和三维模态的通用闭合。由某一算例改善不能推出跨几何适用。

\subsection{局地快率能否预测全局完成时间和尺寸效应}
\textbf{对应文章：}\citet{Birn2001}（2001，Birn2001）；
\citet{L241}（2017，L241，
\emph{The role of guide field in magnetic reconnection driven by island coalescence}）；
\citet{Stanier2013}（2013，Stanier2013）。

\textbf{原文证据与缺口：}Birn2001针对已启动单片；L241图1--10中Hall与PIC的岛合并速率、压力阻力和尺寸依赖在弱导场下显著不同，强导场时某些率接近但拓扑和层宽仍不同；Stanier2013中磁通堆积、反弹及岛是否排出受装置几何控制。
\textbf{本报告综合判断：}局地快通量处理能力不足以唯一给出整个磁岛系统完成合并的时间。需要相同初边值、统一上游归一化和累计磁通诊断；这些论文没有共同建立跨系统的普适完成时间定律。

\subsection{有限三维X线的受抑长度是否普适，Hall模型能否独立复现}
\textbf{对应文章：}\citet{F659}（2019，F659，
\emph{Three-Dimensional Magnetic Reconnection With a Spatially Confined X-Line Extent}）；
\citet{F792}（2020，F792，
\emph{Scaling of Magnetic Reconnection With a Limited X-Line Extent}）。

\textbf{原文缺口：}F659图2--9及F792图2--5给出活跃/受抑区与整体率下降；两篇都是PIC研究，电子非对角压力仍支撑破冻结。约$10d_i$受抑尺度来自特定反平行序列，拟合参数、短层瞬态、导场及$\beta$推广仍有限制。
\textbf{本报告综合判断：}尚不能把该尺度写成标准Hall MHD的普适定律。有限长度与压力闭合是否相互改变，是由这两篇及F611共同引出的研究设计，并非三篇已经证实的结论。\citep{F611}

\subsection{离子加速怎样变成不可逆热化，能量怎样分给两物种}
\textbf{对应文章：}\citet{Yoo2013}（2013，Yoo2013，
\emph{Observation of Ion Acceleration and Heating during Collisionless Magnetic Reconnection in a Laboratory Plasma}）；
\citet{L272}（2021，L272）；
\citet{E512}（2011，E512）。

\textbf{原文缺口：}Yoo2013图2--4确立电势、加速与下游升温的空间联系，却保留碰撞、再磁化、波粒及边界等解释。L272图1--7提出电场梯度失磁化途径，但给定场能量展宽和PIC高温关联不能单独唯一确认所有条件下的热化机制；E512在强导场、低$\beta$排序下建立相混合与弱碰撞熵产链，范围又不同。新增核读的\citet{Drake2009Pickup}反平行拾取模型给中心矩温升标度，但观测系数偏低；其第16段与Yoon2023原文Discussion部分对电场净功和拾取主导程度的归因有分歧。\textbf{本报告判断：}需同轨道逐分量电功预算，不能以末速度、温度或Hall场图单独裁决。本轮修复的是这项比较在报告中缺失，并未宣称解决该研究分歧。
\textbf{11月12日专题补充：}双温Hall并非完全不能研究加热。\citet{L223}（2016）图3已给出两物种不同热点，\citet{Kaptanoglu2020}以三维双温Hall比较HIT-SI实验，\citet{TenBarge2019MMS}（2019）图1已输出两物种热能预算；但分别受超电阻/输运闭合、未取得的勘误及压力热流闭合限制。\citet{Tanabe2015MAST}的热点位置为前一类模拟提供实验约束，尚未唯一验证局地热化机制。

\textbf{比例必须区分：}\citet{F674}正式版图7的约2:1是离子/电子电场沉积率；同图热能变化约1:1，且含储存与边界焓流。\citet{Yoo2017Heating}非对称MRX沉积更接近；\citet{Yoon2023PhaseSpace}仍明确留下精确离子能量份额和有序/随机增能占比。\textbf{本报告综合判断：}不能用Hall总功为零否定物种增能，也不能用局地正电功替代不可逆热化。条件性两温预算已经可算；普适不可逆加热比仍需动理学信息或经独立检验的闭合，并分别排查模型缺失、数值误差、输运边界和诊断不足。独立专题源与PDF位于项目根目录\texttt{specific problems/离子加热问题/}。

\subsection{自然触发阈值与三维后期状态能否由同一模型预测}
\textbf{对应文章：}\citet{Birn2001}（2001，Birn2001）；
\citet{Cassak2005}（2005，Cassak2005）；
\citet{Cassak2010}（2010，Cassak2010）；
\citet{Shepherd2012}（2012，Shepherd2012）。

\textbf{原文缺口：}GEM基准跳过线性触发并明确保留三维问题；双稳态研究的触发与耗散、片层历史相关，简化模型未闭合层长和变化的上游场；Shepherd2012只在斜模明显相互作用前拟合传播，后期不再满足孤立X线图像。
\textbf{本报告综合判断：}这些研究各解决一个阶段，尚不能直接拼成从任意初始电流片到自然启动、三维扩展及饱和的无参数预测。第3章的启动线与第三维线应联合检验，而不能把阶段性公式无条件串接。


\paragraph{如何回到详细学习。}
先抄写第\ref{sec:full-five-moment}节的完整五矩系统，再逐步完成第\ref{sec:twofluid-to-hall}节的Hall约化；接着按本章两套时间步流程理解第2章实现，最后用实验表选择需要比较的诊断量。这样，核心方程、删项、代码和实验就能沿同一条路线对应起来。
